\subsection{The main collar question: the spaces that must be distinguished}

We work in a fixed o-minimal expansion of a real closed field $K$.
In this subsection compact means closed and bounded and every map and family
is definable. A family over a finite parameter polyhedron $P$ is continuous
as a map of total spaces and has a continuous inverse on its image.
Its collar interior is required to be a neighborhood of $B\times P$ in
$X\times P$. These requirements are part of the meaning of a family;
pointwise collarhood alone is not used to substitute for them.

A \emph{strict collar embedding} is an embedding
$c:B\times[0,1]\hookrightarrow X$ with $c(b,0)=b$, whose image minus
$c(B\times\{1\})$ is open in $X$, and which extends to a collar on
$B\times[0,2]$. Equivalently, its outer boundary is collared from the
exterior, after matching the two boundary parametrizations. An
\emph{unparametrized strict collar} is $(C,\dot C)$ with such a
parametrization existing but not chosen. A rigidification chooses $c$.
These are the three distinct objects of 119/148. A \emph{bordant pair}
will mean a closed definable pair $(X,B)$ with a finite definable local
collar atlas. No local-manifold hypothesis is part of this definition.
All statements below about general pairs are made for compact $X,B$;
this is the explicit parenthetical hypothesis in the original 119/148
heading, checked on the scan. A claim for arbitrary noncompact pairs
would be an additional assertion.

The simplicial space of a choice has as its $q$-simplices the indicated
families over $\Delta^q_K$. We say it is contractible when its geometric
realization is contractible. In particular, this is not the assertion
that the set of all definable maps is a definable finite-dimensional set.

\begin{saresult}{A.101}{Proved}{Fibrewise collar automorphisms and cone automorphisms}
For every compact definable $B$, the simplicial group of automorphisms
\[
 (b,t)\longmapsto(b,h(b,t))\quad\text{of }B\times[0,1]
\]
is contractible when $h(b,0)=0$ and $h(b,1)=1$. For every compact
definable $Y$, the simplicial group of homeomorphisms of the closed
cone $CY$ fixing both its base $Y$ and its vertex $o$ is contractible.
These prove the fibrewise form of 119/148 I c) and 119/149 II c),
respectively; they do not prove I c) for arbitrary automorphisms.
\end{saresult}
\begin{proof}
Every $h(b,-)$ is strictly increasing. The formula
$h_s(b,t)=(1-s)h(b,t)+st$ preserves strict increase, endpoints,
continuity and definability, and is stationary at the identity. It
contracts each family simultaneously. The inverse is continuous:
the total map over the compact parameter and base is a continuous
definable bijection between definably compact sets, hence a homeomorphism.

Realize $CY$ as the cone of a bounded embedding of $Y$, with radius
$r\in[0,1]$. For $g\in\Homeo(CY;Y,o)$ put, for $s>0$,
\[
 A_s(g)(y,r)=
 \begin{cases}
 s\,g(y,r/s),&0\leq r\leq s,\\
 (y,r),&s\leq r\leq1,
 \end{cases}
 \qquad A_0(g)=\id.
\]
Multiplication by $s$ here is radial dilation on the embedded cone.
The two formulas agree at $r=s$ because $g$ fixes the base. The inverse
is $A_s(g^{-1})$. At $s=0$ both image and domain of the nonidentity
part lie in the cone of radius $s$, proving continuity, uniformly in
every compact parameter family. The vertex is fixed for every $s$.
This is the Alexander contraction. For $Y=\varnothing$ we define
$C\varnothing=\{o\}$, and the assertion is immediate.
\end{proof}

\begin{saresult}{A.102}{Proved under stated hypotheses}{Parametrized strict collars are contractible}
Suppose a compact definable pair $(X,B)$ admits a strict collar.
The simplicial space $\operatorname{Col}^{\mathrm{rig}}(X,B)$ of all
strict collar embeddings is contractible. This verifies the collar-space
portion of Lead L1 with the stated family convention. It is distinct from
contractibility of rigidifications of one fixed image.
\end{saresult}
\begin{proof}
We give the parameter-dependent argument, including the map that makes
the comparison. Fix one strict collar $c$. Let $d_p$, $p\in P$, be any
family, where $P$ is a finite compact parameter polyhedron. Uniform
definable continuity at $B\times\{0\}\times P$ gives
$\epsilon\in K$, $0<\epsilon<1$, such that
\[
 d_p(B\times[0,\epsilon])\subset c(B\times[0,1))
 \quad\text{for every }p.
\]
Indeed the inverse image of the right-hand open set contains the
definably compact bottom face; its complement has positive distance
from that face. Put $k_p=c^{-1}d_{p,\epsilon}$, where
$d_{p,\epsilon}(b,t)=d_p(b,\epsilon t)$. Thus $k_p$ is a collar
embedding into $B\times[0,1]$, with image disjoint from the top.

Here is the finite formula behind Armstrong's collar comparison
\cite[Theorem 2, pp.~123--124]{Armstrong1970}. For $0\leq s\leq1$,
let $F_s:B\times[0,1]\to B\times[-s,1]$ send
$(b,t)$ to $(b,(1+s)t-s)$. Define $G_{p,s}$ on the latter cylinder by
\[
\begin{aligned}
 G_{p,s}(b,u)&=k_p\left(b,\frac{u+s}{1+s}\right)
 &&(-s\leq u\leq0),\\
 G_{p,s}(k_p(b,v))&=k_p\left(b,\frac{v+s}{1+s}\right)
 &&(0\leq v\leq1),
\end{aligned}
\]
and as the identity outside $k_p(B\times[0,1])$ in the nonnegative
cylinder. The formulas agree on their common boundary. Their images
partition the small collar into radii $[0,s/(1+s)]$ and
$[s/(1+s),1]$, so $G_{p,s}$ is a homeomorphism. Consequently
$K_{p,s}=G_{p,s}F_s$ is a cylinder homeomorphism fixing both ends,
$K_{p,0}=\id$, and
\[
 K_{p,1}(b,t)=k_p(b,t)\quad(0\leq t\leq1/2).
\]
All these maps and inverses are definable continuous families. Glue
$cK_{p,s}c^{-1}$ to the identity on $X\setminus c(B\times[0,1))$.
The top boundary is fixed, so this is an ambient family of homeomorphisms.
It carries $c_{1/2}$ to $d_{p,\epsilon/2}$. Concatenate the shrinking
path from $c$ to $c_{1/2}$, this ambient path, and the expansion from
$d_{p,\epsilon/2}$ to $d_p$. Every intermediate embedding is still a
strict collar, by its extension or by the ambient homeomorphism.
Thus every finite parameter family is homotopic to a constant one.

A finite simplicial sphere of choices is such a family, by gluing its
finitely many defining formulas. Hence all homotopy groups of the
realization vanish and it is nonempty. A simplicial realization has
CW homotopy type, so Whitehead's theorem makes it contractible.
The proof also covers $B=\varnothing$, when the choice is unique.
\end{proof}

\begin{saresult}{A.103}{Known}{A fully explicit nonzero concordance group}
Let $B=(S^1)^d$, $d\geq6$, with its finite semialgebraic PL model.
For $\mathrm{CAT}=\mathrm{PL}$ or $\mathrm{TOP}$, put
\[
 P_{\mathrm{CAT}}(B)=
 \Homeo_{\mathrm{CAT}}(B\times[0,1]\,;B\times\{0\}).
\]
Then
\[
 \pi_0P_{\mathrm{CAT}}(B)
 \cong \F_2[t_1^{\pm1},\ldots,t_d^{\pm1}]/\F_2\{1\}
 =\bigoplus_{v\in\Z^d\setminus\{0\}}\F_2.
\]
The same component computation holds for definable concordances of this
model in every o-minimal expansion of a real closed field. In particular
we may take $B=(S^1)^6\subset K^{12}$, given by the six equations
$x_i^2+y_i^2=1$.
\end{saresult}
\begin{proof}[Derivation from the opened results]
We first correct a potentially consequential source issue. The
general split formula in Hatcher's Theorem~3.1 cannot be used
unchanged. The corrected smooth sequence is
\[
 \operatorname{Wh}_3(\pi)\longrightarrow
 \operatorname{Wh}_1^+(\pi;\Z/2)\longrightarrow
 \pi_0P_{\rm Diff}(B)\longrightarrow
 \operatorname{Wh}_2(\pi)\longrightarrow0,
 \qquad \pi=\pi_1B,
\]
when $\dim B\geq5$ and $\pi_2B=0$
\cite[Theorem~7.1, p.~23; Remark~7.3, p.~26]{Igusa2025}.
In the definition on p.~2 of that source,
$\operatorname{Wh}_3(\pi)$ is $K_3(\Z[\pi])$ modulo the constant
$K_3(\Z)$ and the framed assembly image from
$\Omega^{\rm fr}_3(B\pi)$. We verify the missing term for the torus.

The Laurent fundamental theorem for regular rings, including its
splitting by multiplication by the unit classes $[t_i]$, gives
\[
\begin{split}
K_3(\Z[\Z^d])\cong{}&K_3(\Z)
 \oplus\bigoplus_iK_2(\Z)[t_i]
 \oplus\bigoplus_{i<j}K_1(\Z)[t_i][t_j]\\
 &\oplus\bigoplus_{i<j<k}K_0(\Z)[t_i][t_j][t_k].
\end{split}
\]
There are no Nil terms, by regularity
\cite[Theorem~6.3, Corollary~6.3.1, p.~40; Theorem~8.2, p.~61]{WeibelV}.
The unit from stable homotopy to algebraic $K$-theory is
multiplicative and sends $\eta$ to $[-1]$
\cite[Exercise~1.12, pp.~15--16; Corollary~1.20, p.~13]{WeibelIV}.
Hence $\eta^2$ maps to $\{-1,-1\}$, the generator of
$K_2(\Z)=\Z/2$ \cite[Example~5.2.2, p.~36]{WeibelIII}.
The classes $\eta^2$ times each coordinate circle, $\eta$ times
each coordinate two-torus, and the framed coordinate three-tori
therefore hit every nonconstant summand. The constant summand
is quotiented out separately. Thus
$\operatorname{Wh}_3(\Z^d)=0$. In degree two the identical
Laurent decomposition has only the constant $K_2$, the unit
$K_1[t_i]$ summands and $K_0[t_i][t_j]$ summands; their
monomial generators are precisely the defining relations in
$\operatorname{Wh}_2$, so that group is also zero.
Because $\Z^d$ acts trivially by conjugation, the middle
Whitehead term is $\F_2[\Z^d]/\F_2[0]$. Exactness proves the
smooth formula without any splitting assertion.

For the category comparison, Burghelea--Lashof's
\cite[Theorem~6.1, p.~41; proof pp.~42--43]{BurgheleaLashof1974II}
states directly that smooth-to-PL concordances induce an
isomorphism on components for a compact connected smooth
manifold. Its proof identifies the relative smoothing space
with a section space and proves its components and fundamental
group vanish. This is an unstable comparison, exactly the map
we need for the smooth torus with its differentiable
triangulation. Theorem~6.2 of the same source supplies the
PL--TOP comparison in the safe range $d\geq6$ adopted here.
Both fix the bottom and the side boundary; the latter is empty
for the torus. Thus no stability inference is needed.
The torus formula is also explicitly recorded in
\cite[\S6.2, p.~34, equation~(6.12), Lemmas~6.16--6.17]{CourtePorcelli2025};
we do not use those stable statements as a substitute for the
unstable comparison just checked.
The nonconstant monomial $t_1$ is nonzero in the displayed
additive quotient, proving nontriviality.

For the definable comparison use the relative homeomorphism and isotopy
arguments in Shiota \cite[Theorem 2.1 and its complement; Lemma 3.3,
proof of injectivity, pp.~44--46]{Shiota2010}. They give two precise
facts for a compact polyhedral pair: every definable self-homeomorphism
fixed on a subpolyhedron is definably isotopic, relative to it, to a PL
one; and two such PL homeomorphisms joined by a definable isotopy are
joined by a PL isotopy. Apply them to
$(B\times[0,1],B\times\{0\})$, keeping the other boundary component
as a distinguished subpolyhedron. Shiota's Lemma 1.2 identifies the
PL isotopy classes over different ordered fields. Thus inclusion of PL
concordances induces a bijection on components, compatible with
composition. We are using a component theorem, not asserting from
these locators alone a comparison of all higher homotopy groups.
\end{proof}

\begin{saresult}{A.113}{Known}{The general component comparison and the PL--TOP comparison}
Let $B$ be any closed definable manifold and choose its finite PL
model, with the closed-manifold convention including definable
compactness. Inclusion gives a natural isomorphism
\[
 \pi_0P_{\rm PL}(B)\xrightarrow{\sim}\pi_0P_{\rm def}(B).
\]
These component groups are independent of the ordered field used
to realize the finite PL model. Over the real field, if every
component of $B$ has dimension at least five, the map of
simplicial groups $P_{\rm PL}(B)\to P_{\rm TOP}(B)$ is a weak
equivalence. Thus for general $K$ the definable component group
also agrees with the TOP component group of a chosen real
realization of the finite PL model in that range. We do not claim
a canonical map from a non-Archimedean realization to the real
TOP space.
The last statement does not assert a weak equivalence with the
whole definable concordance group.
\end{saresult}
\begin{proof}[Derivation from the opened results]
The two relative facts applied in the proof of A.103 did not use
the torus. Apply the same relative homeomorphism straightening
and relative isotopy theorem of \cite[Theorem 2.1 with complement;
Lemma 3.3, pp.~44--46]{Shiota2010} to the compact polyhedron
$B\times I$, fixing $B\times\{0\}$ pointwise and preserving the
other boundary component. The first theorem replaces every
definable concordance by a PL concordance through such maps,
proving surjectivity. The relative injectivity argument takes an
isotopy between two PL maps, triangulates the images of finitely
many simplex cylinders while preserving the single isotopy
coordinate, and uses the level-preserving Alexander extension
on each simplex cylinder, fixed on its already treated boundary.
It gives a PL isotopy with the same two endpoints and the same
fixed subpolyhedron. This proves injectivity. Composition is
unchanged by the inclusion, so the bijection is a group
isomorphism. Shiota's Lemma 1.2 transports these PL isotopy
classes between real closed fields.

Burghelea--Lashof's Theorem 6.2, printed pp.~41--43,
identifies the homotopy groups of PL and TOP concordances for
compact connected PL manifolds in dimension at least five
\cite{BurgheleaLashof1974II}; the exact range is also stated in
the opened \cite[Remark 4.7, p.~35]{GoodwillieKrannichKupers2023}.
Their concordances fix $\partial B\times I\cup B\times\{0\}$,
which is exactly our condition because $\partial B$ is empty.
The PL simplicial model and the singular complex of the TOP
group are the models used in the theorem. A finite disjoint
union reduces to a finite product, since fixing the bottom
pointwise prevents permutation of its components. This verifies
every hypothesis and gives the asserted weak equivalence.
\end{proof}

\begin{saresult}{A.104}{Refuted}{I c) and the proposed collapse equivalence}
For $B=(S^1)^6$, the automorphism space in 119/148 I c) is not
contractible. Moreover the boundary-collapse functor from bordant pairs
to pointed spaces is not a Dwyer--Kan equivalence in the stipulated
definable simplicial categories. Its map on automorphism spaces at
$(B\times[0,1],B\times\{0\})$ has homotopy fibre equivalent to
$P_{\mathrm{def}}(B)$. Therefore the full cascade of 119/150 and Lead
L1 cannot be repaired by rigidifying collars alone.
\end{saresult}
\begin{proof}
The space in I c) is precisely $P_{\mathrm{def}}(B)$: only its lower
boundary is fixed pointwise, and its upper boundary is free.
A.103 proves that its component group is nonzero.

Write $I=[0,1]$ and put
\[
 G=\Homeo_{\mathrm{def}}(B\times I;B\times\{0\}\text{ setwise}),
 \qquad H=\Homeo_{\mathrm{def}}(CB;o).
\]
The upper boundary of the cylinder and the base of the cone are
intrinsic boundaries of their manifold loci. The cone vertex is the
specified point. Hence upper-boundary restriction
$e:G\to\Homeo_{\mathrm{def}}(B)$ and base restriction
$r:H\to\Homeo_{\mathrm{def}}(B)$ are defined. Both have sections,
by constant extension and coning, respectively.
The kernel of $e$ is the group of cylinder homeomorphisms fixed on the
upper boundary; interval reversal identifies it with $P_{\mathrm{def}}(B)$.
Multiplying by the constant extension of $e(g)^{-1}$ gives an actual
product decomposition of the underlying simplicial spaces of $G$.
Similarly the kernel of $r$ consists of cone homeomorphisms fixed on
base and vertex and contracts by A.101. The same decomposition makes
$r$ a homotopy equivalence.

The quotient $q:B\times I\to CB$ collapses the lower boundary.
Every $g\in G$ descends uniquely to $\bar g\in H$, and
\[
 r(\bar g)=e(g).
\]
Thus the homotopy fibre of $G\to H$ is that of $e$, namely
$P_{\mathrm{def}}(B)$. By A.103 it is not contractible. The induced
map on this mapping space is not a weak equivalence, which disproves
the required enriched equivalence.
Equivalently choose any concordance representing $t_1$ in A.103 and
reverse its interval. It is nontrivial in $\pi_0G$, by the split
upper-boundary restriction, whereas its descended cone automorphism
lies in the Alexander-contractible kernel of $r$.
We do not assert that every bottom-fixed concordance descends to that
kernel: its upper endpoint need not be the identity.
\end{proof}

\begin{saresult}{A.105}{Proved}{Rigidified blow-up and forgetting cone rigidification}
Use strict collars and strict conical neighborhoods, and require morphisms
to preserve the chosen product or cone parametrizations. Then
\[
 PBCR\simeq EPCR\simeq EPC.
\]
The first is an equivalence induced by collapsing the lower collar
boundary and blowing up the cone vertex. The second is a weak equivalence
on every mapping space and is essentially surjective. These assertions
hold with any finite number of disjoint labeled boundary pieces or points.
They do not imply equivalence with $EP$.
\end{saresult}
\begin{proof}
For $(X,B,c)$ collapse $B$ and identify the image of $c$ with
$C\dot C$ by its radial parameter. Conversely for $(Z,o,C,\psi)$ put
$E=Z\setminus(C\setminus\dot C)$ and form
\[
 \widetilde Z=E\amalg_{\dot C}(\dot C\times[0,1]),
\]
where $\dot C\times\{1\}$ is attached to $\dot C\subset E$.
The new lower boundary is $\dot C\times\{0\}$ and its displayed
product is the rigidified collar. Strictness is exactly the condition
that the outer face is collared in $E$. Each of these two constructions
followed by the other has the canonical indicated homeomorphism to the
starting object. A morphism respecting radial parameters extends
uniquely across the blown-up lower boundary; a product-preserving
morphism descends uniquely. These statements hold on total spaces of
families and respect simplicial operators. They prove the first
equivalence, including its inverse on all maps.

For the second functor fix any two rigidified cone objects. Given a map
$f$ preserving the unparametrized cone, replace its restriction inside
that cone by
\[
 \psi'\,C(f|_{\dot C})\,\psi^{-1}
\]
and leave it equal to $f$ outside. The formulas agree on $\dot C$.
The discrepancy is a homeomorphism of the source cone fixing base and
vertex. Contract it by A.101 and extend the contraction as the identity
outside the cone. This gives a deformation retraction of mapping spaces
onto the maps preserving the rigidification. Every unrigidified cone has
a rigidification by its definition, proving essential surjectivity.
For finitely many disjoint cones use the formulas separately on each
closed cone; they agree on the complementary region. Permutations of
labels commute with the constructions. Empty boundary pieces correspond
to isolated cone vertices under $C\varnothing=\{o\}$.
\end{proof}

\begin{saresult}{A.106}{Refuted}{Unparametrized collars retain the classifying space of concordances}
Use parametrizable families for the unparametrized collar simplicial set:
a simplex is a strict collar family admitting a parametrization, with no
parametrization retained. If $(X,B)$ admits a collar, then
\[
 \operatorname{Col}(X,B)\simeq BP_{\mathrm{def}}(B).
\]
In particular 119/148 I b) is false for $B=(S^1)^6$:
$\pi_1\operatorname{Col}(X,B)\cong\pi_0P_{\mathrm{def}}(B)\ne0$.
The rigidified assertion A.102 and the fibrewise assertion A.101 are
the precise different assertions that remain true.
\end{saresult}
\begin{proof}
Precomposition makes $E=\operatorname{Col}^{\mathrm{rig}}(X,B)$ a free
simplicial $P_{\mathrm{def}}(B)$-set. Two parametrizations have exactly
the same collar and outer face if and only if their transition map is
in this group. Thus the stated unparametrized space is the degreewise
quotient $E/P_{\mathrm{def}}(B)$.

This is a principal Kan fibration. Here are the needed details for the
family convention. Both the group and the embedding space satisfy horn
filling: glue the compatible maps on a finite simplicial horn and pull
the resulting family back by a PL retraction of the simplex onto the
horn. Such a retraction is obtained by an elementary simplicial collapse.
For a horn lifting problem for the quotient, choose a representative of
the given quotient simplex. Its transition maps to the prescribed horn
lifts form a horn in the group, since the action is free. Fill that
group horn and adjust the representative. This proves the lifting
property. A.102 contracts the total space. The long exact homotopy
sequence and the universal principal-bundle construction therefore
identify the base with the classifying space of the group. The
nonvanishing follows from A.103.
\end{proof}

\begin{saresult}{A.107}{Known}{Local collars glue over a real closed field}
Let $B$ be closed in a definable subset $X$ of an affine space over a
real closed field. If $B$ has a finite definable local collar atlas, then
there is a definable open collar $B\times[0,1)\to X$. If $X$ and $B$
are definably compact, it restricts after rescaling to a strict closed
collar. Thus I a) holds for the compact bordant pairs used here.
\end{saresult}
\begin{proof}[Derivation and definability check]
Brown's finite gluing lemma is \cite[Lemmas 1--4, pp.~333--337]{Brown1962}:
the union of two relatively open collared subsets is collared. Its proof
uses distance functions, minima, finite closed shrinkings of an open
cover, and piecewise affine changes of the interval coordinate. In a
definable subset over a real closed field these operations are definable.
Distances to definable sets exist by definable completeness; a distance
to the complement is positive at every point of an open set. Closed
shrinkings and separated neighborhoods are given by the same distance
inequalities. Brown's spindle modification is the identity below half
a positive height and conjugates the transition map by the linear
collapse above that height. Its inverse has the displayed same finite
piecewise formula, so its continuity proof uses no sequential
completeness or infinite process. Hence Lemma 4 holds verbatim for a
finite definable collar atlas. Induction on its finite size yields an
open collar. In the compact case its restriction to $[0,1/2]$ is a
closed embedded collar, and restriction to $[0,3/4]$ supplies an outer
extension, proving strictness.
\end{proof}

\begin{saresult}{A.108}{Known}{Existence of strict pointed cones}
Every compact definable pointed space $(X,o)$ admits a strict closed
conical neighborhood of $o$, including when $o$ is isolated.
\end{saresult}
\begin{proof}[Derivation]
Take a finite definable triangulation compatible with $\{o\}$, so that
$o$ is a vertex. Its barycentric coordinate $\lambda_o$ is a continuous
PL function, extended by zero off the closed star. For $0<\epsilon<1$,
the set $C=\{\lambda_o\ge\epsilon\}$ is a cone with vertex $o$ on the
level set $\{\lambda_o=\epsilon\}$. Varying the coordinate in
$[\epsilon/2,\epsilon]$ gives an exterior collar of its outer face.
The complement of that outer face in $C$ is open around $o$. If the
vertex is isolated the level set is empty and $C=\{o\}$. The finite
triangulation input is the definable triangulation theorem used in
\cite[\S1]{Shiota2010}; the cone and strictness are given explicitly here.
\end{proof}

\begin{saresult}{A.109}{Proved}{Attaching a collar and attaching a cone are enriched equivalences}
Let $F:PB\to PBCR$ attach $B\times[-1,0]$ to $X$ along
$B\times\{0\}=B$, with new boundary $B\times\{-1\}$ and its
displayed rigid collar. Let $K:PB\to EPCR$ attach $CB$ to $X$
along its base. Both functors are equivalences of the specified
definable simplicial categories, with actual isomorphisms on mapping
spaces. In particular the valid corrected comparison is
\[
 PB\simeq PBCR\simeq EPCR\simeq EPC.
\]
This conclusion does not include $EP$.
\end{saresult}
\begin{proof}
The complement of the attached collar minus its outer face is exactly
the original $X$. A map respecting the displayed rigid collar restricts
to a homeomorphism of $(X,B)$; conversely any such homeomorphism extends
uniquely by the product of its boundary restriction and the interval
identity. These operations are inverse on all parameter families.
The inverse construction on an arbitrary rigid collar object is
\[
 (Z,B,C,c)\longmapsto
 \bigl(Z\setminus(C\setminus\dot C),\dot C\bigr).
\]
Strictness says precisely that this is a bordant pair. The chosen
parametrization identifies the collar reattached to it with the
original collar, giving a natural isomorphism to the starting object.
Thus $F$ is fully faithful and essentially surjective, with its stated
inverse.

For $K$ use the same exterior and attach a cone. A rigid cone map is
uniquely the cone of its base map, and its exterior restriction is a
bordant-pair map. Hence the identical argument proves the assertion
for $K$. The other equivalences in the display are A.105. Our pointed
quotient convention is $X/B=(X\amalg\{o\})/(B\sim o)$; when
$B=\varnothing$ this is $X\amalg\{o\}$. Thus empty boundary pieces
correspond to isolated marked points and all inverse constructions
remain valid. Finitely many disjoint labeled pieces are treated by
the same formula on each piece; the formulas commute with permutations
of the labels.
\end{proof}

\begin{saresult}{A.110}{Refuted}{The naturality qualification for the two displayed $2$-cells}
On $PB$ with all homeomorphisms, neither the collar-absorption
isomorphism $UF\Rightarrow\id$ nor the cone-absorption isomorphism
$VK\Rightarrow Q$ is a strict natural isomorphism. Here $U$ forgets
the attached collar, $V$ forgets the attached cone, and $Q$ collapses
the original boundary. They become strict natural isomorphisms after
passing to pairs with a chosen extended collar and maps preserving its
product coordinates. This is the role of the manuscript's marginal
instruction to precompose with $PBR$ at 119/150.
\end{saresult}
\begin{proof}
Take $X=[0,1]$ and $B=\{0\}$. Every increasing definable
homeomorphism of this interval fixes the newly attached collar or cone
pointwise, because it fixes $B$. The common fixed set of these
homeomorphisms on the original interval is exactly $\{0,1\}$:
each interior point is moved by an increasing piecewise affine
homeomorphism. Naturality would force the image of the entire attached
interval under either proposed isomorphism to lie in this two-point
set. No homeomorphism has this property. This proves both failures.

Now choose $c:B\times[0,2]\hookrightarrow X$. In the union of its
image with the attached collar use the coordinate $t\in[-1,2]$ and
send it to $2(t+1)/3\in[0,2]$, keeping $b$ fixed; use the identity
outside this extended collar. The formulas agree at $t=2$, give a
homeomorphism, and send the new boundary to the old one. They commute
with every map preserving the chosen product coordinates. This is
$\alpha$ on that category. Quotienting its source and target by their
respective boundary pieces gives $\beta$, because the source quotient
is exactly the attached-cone construction. This also proves the
compatibility of the two $2$-cells. The piecewise construction applies
independently to each of finitely many disjoint chosen collars.
\end{proof}

\begin{saresult}{A.111}{Refuted}{The strict-cone choice obstruction in the stipulated categorical model}
For a simplicial category of homeomorphisms, define its derived space
of choices over an object to be the homotopy fibre of the map on
classifying spaces of its maximal simplicial groupoids. For
$X=C B$ with $B=(S^1)^6$, the component of the standard strict cone
in the derived space of strict conical neighborhoods of the vertex
has homotopy type
\[
 BP_{\mathrm{def}}(B).
\]
Consequently II b) is false with this categorical interpretation of
choice, and remains false for rigidified strict-cone choices. The
space of rigidifications of a \emph{fixed} cone is contractible, as
II c) asserts. These three assertions concern different spaces.
\end{saresult}
\begin{proof}
Give $CB$ radial interval $[0,2]$ and choose the smaller cone
$C_1B$ with radii $[0,1]$. Its exterior is $B\times[1,2]$.
Let $H$ be the simplicial group of homeomorphisms of $CB$ fixing
the vertex, and let $J\subset H$ preserve the smaller cone.
For this torus link the outer base is the intrinsic manifold boundary
of the cone's manifold locus, so every member of $H$ preserves it.
Restriction to the outer base gives split homomorphisms
\[
 H\longrightarrow\Homeo(B),\qquad
 J\longrightarrow\Homeo(B).
\]
The first kernel is Alexander-contractible by A.101. Write $J_0$
for the second kernel. Restriction of $J_0$ to the shell gives a
split homomorphism
\[
 J_0\longrightarrow
 \Homeo(B\times[1,2];B\times\{2\}),
\]
whose section cones the map induced on the inner face. Its kernel
is the group of maps of $C_1B$ fixed on its base and vertex, again
Alexander-contractible. The target is $P_{\mathrm{def}}(B)$ after
reversing and rescaling the interval. These split decompositions are
decompositions of simplicial spaces; hence $J_0\simeq P_{\mathrm{def}}(B)$.
It follows, by applying classifying spaces to the split group maps,
that the homotopy fibre of $BJ\to BH$ is $BJ_0\simeq BP_{\mathrm{def}}(B)$.
The component of an object in the classifying space of a simplicial
groupoid is the classifying space of its automorphism group. Therefore
this is exactly the indicated component of the derived choice space.
Surjectivity of $\pi_0J\to\pi_0H$ makes this fibre connected. Other
components of the cone-object category cannot remove its nonzero
fundamental group.

A.103 gives $\pi_1BP_{\mathrm{def}}(B)=\pi_0P_{\mathrm{def}}(B)\ne0$.
Finally A.105 identifies the rigidified and unrigidified cone categories
over the pointed-space category by weak equivalences on mapping spaces,
so their derived choice spaces agree. This argument does not invoke an
unproved parametrized isotopy-extension theorem to identify a different,
pointwise-defined topology on a set of cone neighborhoods with this
derived space.
\end{proof}

\begin{saresult}{A.112}{Proved}{The eleven arrows on 119/150}
For the strict collar and strict cone conventions above, five of the
eleven arrows are equivalences of the definable simplicial categories.
The other six fail in the torus example of A.103. The following table
distinguishes the two copies of $PB$ in the drawing by position.
\begin{center}
\begin{tabular}{p{0.58\linewidth}p{0.30\linewidth}}
\toprule
Arrow & Outcome\\
\midrule
Top $PB\to PBCR$, attach a collar & equivalence (A.109)\\
Top $PB\to EPCR$, attach a cone & equivalence (A.109)\\
Top $PB\to$ bottom $PB$, identity & equivalence\\
Top $PB\to EP$, collapse & fails (A.104)\\
$PBCR\to EPCR$, collapse & equivalence (A.105)\\
$PBCR\to PBC$, forget product coordinates & fails\\
$EPCR\to EPC$, forget cone coordinates & equivalence (A.105)\\
$PBC\to EPC$, collapse & fails\\
$PBC\to PB$, forget the collar & fails\\
$EPC\to EP$, forget the cone & fails (A.111)\\
Bottom $PB\to EP$, collapse & fails (A.104)\\
\bottomrule
\end{tabular}
\end{center}
Here ``fails'' means that the functor is not an equivalence in general;
we do not assert failure at every bordant pair with that boundary.
\end{saresult}
\begin{proof}
Only the three remaining collar arrows require verification. Put
$X=B\times[0,2]$, with boundary at $0$ and chosen collar $[0,1]$.
Let $J_c$ be its collar-preserving automorphism group. Restriction to
the upper half gives a split map
\[
 J_c\longrightarrow G_+=\Homeo(B\times[1,2];B\times\{1\}
 \text{ setwise}).
\]
Its section is constant extension of the inner boundary map across
the lower half. Its kernel is the lower-half group fixed on its top,
which is $P_{\mathrm{def}}(B)$ after reversal. The automorphism group
with rigid collar is exactly this section $G_+$. Since the underlying
space of $J_c$ is a product of $G_+$ with this kernel, the section
fails to induce a surjection on components. This disproves
$PBCR\to PBC$ at the displayed object.

The inclusion of the rigid-collar automorphisms into all automorphisms
of $(X,B\times\{0\})$ is a homotopy equivalence in this example.
To see it explicitly, write an automorphism as the constant extension
of its lower boundary map times a bottom-fixed $p$. For $0\le s\le1$
compress $p$ into $[s,2]$: use the identity below $s$ and
$D_s pD_s^{-1}$ above $s$, where
$D_s(b,t)=(b,s+(2-s)t/2)$. The formulas agree at $s$ because $p$
fixes the bottom, and vary continuously and definably, including
$s=0$. At $s=1$ the result preserves the rigid collar. If it already
preserved that collar, this entire homotopy stays within the same
subgroup. It therefore proves the claimed homotopy equivalence.
If forgetting the unrigidified collar were also a weak equivalence
on this mapping space, the two-out-of-three property would contradict
the failure of the preceding section. Thus $PBC\to PB$ fails.

After collapse, restriction to the unchanged upper shell identifies
the cone-preserving group with $G_+$ up to the Alexander-contractible
kernel of maps of the small cone fixed on its base and vertex.
The map from $J_c$ commutes with this restriction. Its lower-half
kernel $P_{\mathrm{def}}(B)$ maps into that contractible cone kernel.
The same split product decompositions show that its homotopy fibre
is $P_{\mathrm{def}}(B)$. Hence $PBC\to EPC$ fails too. The other
entries are the stated earlier results. Taking one nontrivial boundary
piece and leaving all other labels fixed gives the same failures in
every finite multipair variant in which that piece is permitted.
\end{proof}

